\section{Experimental Validation}
We conduct a comprehensive evaluation of the RePAIR framework and the proposed STAMP unlearning method to answer three key research questions. \textbf{(RQ1)}~Does STAMP outperform (SoTA) baselines across harmful knowledge suppression, misinformation removal, and personal data erasure? \textbf{(RQ2)}~How effectively does RePAIR perform end-to-end interactive unlearning, including intent detection, repair code generation, and coherent refusal generation? \textbf{(RQ3)}~Does the pipeline behave correctly end to end, and does the resulting forgetting survive reformulation of the query? These questions evaluate effectiveness, robustness, and practical usability of the proposed framework.\\
\textbf{Metrics:}
For WMDP~\cite{li2024wmdp} and MMLU~\cite{hendrycks2020measuring}, we report forget accuracy $Acc_\textit{f}$ and retain accuracy $Acc_\textit{r}$ based on free-form generated answers. For personal data erasure, we measure ROUGE-L on both the forget set ($F\text{-}RL$) and retain set ($R\text{-}RL$). Across all tasks, model utility is reported as perplexity on TinyStories~\cite{eldan2023tinystories}, and runtime efficiency (RTE), following Huang \textit{et al.}~\cite{huang2025unified}, is reported in minutes. Ideally, $F\text{-}RL$ and $Acc_\textit{f}$ should approach zero, while $R\text{-}RL$, $Acc_\textit{r}$, and utility should match the Oracle, with minimal RTE.

\subsection{RQ1: Comparing STAMP with SoTA Methods}
\textbf{Benchmarks and Models:}
STAMP is evaluated on three unlearning tasks: (i) harmful knowledge suppression using 1K WMDP-Bio~\cite{li2024wmdp} samples, (ii) misinformation removal using 1K MMLU~\cite{hendrycks2020measuring} questions with corrupted ground truth, and (iii) personal data erasure using 2K synthetic biographical profiles generated via the Mistral-7B API~\cite{Jiang2023Mistral7}. Each dataset is split equally into $\mathcal{D}_f$ and the full retain set $\mathcal{D}_r^{\textit{full}}$; following Section~\ref{sec:PS}, the retain buffer $\mathcal{D}_r$ used by all methods then comprises 10\% of $\mathcal{D}_r^{\textit{full}}$, which the retain-buffer ablation in the supplementary material (Section~C.3) shows is sufficient. A shared $\mathcal{D}_{\textit{ref}}$ of 200 refusal prompts is used for steering vector computation (Section~\ref{sec:stamp}).

For both WMDP and MMLU, we use reduced subsets and train the model to generate free-form answers rather than selecting from MCQ options. Two properties of the multiple-choice format make it unsuitable here. First, its floor is chance rather than zero, which is 25\% for four-way questions, and Llama-3-8B scores ${\sim}67$ on five-shot MCQ MMLU~\cite{grattafiori2024llama}, so the entire measurable effect of unlearning spans roughly forty points and a perfectly forgotten model is indistinguishable from one that is guessing. Second, a refusal cannot be expressed as an option, so the behaviour RePAIR is designed to produce has no place in the answer space, and grafting a refusal onto a fixed set of choices would put the model out of format and leave the resulting response ambiguous to score. Free-form generation removes both problems, since $Acc_f$ can fall to zero and the refusal is directly observable. Reported accuracies are consequently not directly comparable to standard MCQ-based results~\cite{grattafiori2024llama}; in particular, the Base scores in Table~\ref{tab:STAMP-Comparison} (75.30--86.30) exceed the MCQ figures because the model is credited for the content of its generated answer on the reduced subsets rather than for selecting among adversarial distractors. Llama-3-8B~\cite{grattafiori2024llama} serves as $\mathcal{M}_{\textit{patient}}$, Mistral-7B~\cite{Jiang2023Mistral7} as $\mathcal{M}_{\textit{watchdog}}$ for intent classification and forget-pair extraction, and Qwen2.5-Coder-7B-Instruct~\cite{hui2024qwen2} as $\mathcal{M}_{\textit{surgeon}}$ for repair code generation. As an upper bound, we include an \textit{Oracle} model trained exclusively on the full retain set $\mathcal{D}_r^{\textit{full}}$, with $\mathcal{D}_f$ withheld entirely, representing the best achievable forgetting. Six baselines are compared: GA~\cite{yao2024large}, NPO~\cite{zhang2024negative}, RMU~\cite{li2024wmdp}, FLAT~\cite{wang2024llm}, WGA~\cite{wang2025rethinking}, and ASU~\cite{zade2026attention}. Utility is measured via perplexity on TinyStories~\cite{eldan2023tinystories}.
\begin{table}[h!]
\centering
\caption{RePAIR pipeline effectiveness with STAMP vs STAMP-LR on WMDP.}
\label{tab:decomposition}
\renewcommand{\arraystretch}{1.25}
\begin{tabularx}{\columnwidth}{l|>{\centering\arraybackslash}X>{\centering\arraybackslash}X}
\Xhline{1.15pt}
\textbf{Metric} & \textbf{STAMP} & \textbf{STAMP-LR} \\
\Xhline{1.15pt}
Is Valid Python Code (\%)       & 97.23 & 96.27 \\
\rowcolor{gray!15}
User Request Detected (\%)      & 96.30 & 97.50 \\
User Request Satisfied (\%)     & 98.90 & 97.70 \\
\rowcolor{gray!15}
"I don't Know" Rate (\%)                  & 98.27 & 96.27 \\
Turnaround Time (min)          & 9.36  & 6.50  \\
\Xhline{1.15pt}
\end{tabularx}
\end{table}

\medskip
\noindent
\textbf{Results Discussion:}
Table~\ref{tab:STAMP-Comparison} demonstrates STAMP's effectiveness over SoTA baselines across all three tasks. \colorbox{green!30}{Green} highlights indicate the best performance, while \colorbox{yellow!40}{yellow} indicates second-best performance.
\textit{Forgetting and retention:}
Most methods drive $Acc_\textit{f}$ close to zero, but WGA and FLAT retain a clear residual: WGA leaves $Acc_\textit{f}$ at 2.10 on harmful knowledge and 2.47 on misinformation, and records the highest $F\text{-}RL$ of 0.45 on personal data erasure, while FLAT leaves $Acc_\textit{f}$ at 1.30 on misinformation and $F\text{-}RL$ at 0.33. Only STAMP and STAMP-LR reach $Acc_\textit{f} = 0.00$ and $F\text{-}RL = 0.00$ on all three tasks. On retention, RMU maintains the highest $Acc_\textit{r}$ among baselines (74.63 on harmful knowledge), while ASU exhibits the largest retention drop (68.39), likely due to over-smoothing of attention. Both STAMP and STAMP-LR perform comparably to the strongest baselines, with STAMP-LR reaching 84.47 $Acc_\textit{r}$ on misinformation removal, closely matching the Oracle (85.30).
\textit{Utility preservation:}
As anticipated from Section~\ref{sec:related}, GA-based methods (GA~\cite{yao2024large}, NPO~\cite{zhang2024negative}, WGA~\cite{wang2025rethinking}) suffer significant utility degradation, with perplexity rising to 9.27--11.99 on TinyStories. In contrast, FLAT~\cite{wang2024llm}, ASU~\cite{zade2026attention}, STAMP, and STAMP-LR remain between 6.02 and 8.36, close to the Oracle (5.01--6.10) and substantially better than training-based baselines.

\noindent
\textbf{Runtime efficiency:}
All baselines are trained for two epochs, requiring approximately 12 minutes for harmful knowledge removal, 6 minutes for misinformation removal, and 10 minutes for personal data erasure. Being training-free, STAMP reduces this to 7.13, 4.25, and 6.48 minutes, respectively, while STAMP-LR further improves to 4.25, 2.57, and 4.01 minutes, achieving up to $\sim$3$\times$ speedup over training-based methods.

\subsection{RQ2: RePAIR Framework Effectiveness}

We evaluate the full RePAIR pipeline, where intent detection is performed by $\mathcal{M}_{\textit{watchdog}}$, repair code generation by $\mathcal{M}_{\textit{surgeon}}$, and unlearning execution by $\mathcal{M}_{\textit{patient}}$. Experiments use the WMDP dataset~\cite{li2024wmdp} under the same setup as Table~\ref{tab:STAMP-Comparison}, with Llama-3-8B as $\mathcal{M}_{\textit{patient}}$ and the 10\% retain buffer: each of the 500 forget samples is issued as a separate unlearning dialogue, and all rates are averaged over these 500 requests. We report five metrics: valid code rate, request detection, request satisfaction, IDK rate, and turnaround time in Table~\ref{tab:decomposition}. All metrics except turnaround time are evaluated using Mistral-7B~\cite{Jiang2023Mistral7}, while the user role is simulated via a separate Mistral-7B API instance. Results compare STAMP and STAMP-LR.

Both variants achieve above 96\% across all metrics. The high valid code rate is driven by Qwen2.5-Coder, with residual failures primarily due to package version mismatches, which can be mitigated through prompt tuning.

Turnaround times 9.36 and 6.50 minutes exceed the RTE reported in Table~\ref{tab:STAMP-Comparison} due to the additional overhead of multi-model orchestration, including $\mathcal{M}_{\textit{watchdog}}$ and $\mathcal{M}_{\textit{surgeon}}$, alongside the unlearning execution.

\subsection{RQ3: Pipeline in Action and Robustness to Reformulation}
The supplementary material (Section~D) traces a single personal data erasure request end to end. $\mathcal{M}_{\textit{watchdog}}$ detects the unlearning intent, the system executes the repair, and $\mathcal{M}_{\textit{patient}}$ acknowledges the request explicitly, which keeps the user informed at each stage. Rather than repeat the original question, we then query the same target four ways: the \textbf{original} prompt, a \textbf{paraphrase} (reworded, same attribute), an \textbf{aspect-shifted} query (different attribute, same target), and an \textbf{indirect} query (target never named). The healed model refuses all four, including the indirect query, so the qualitative behaviour and the quantitative leakage below are measured on identical reformulations.

To quantify this, we use Mistral-7B to generate the same three reformulations for 200 forget items from the personal data erasure task, querying the healed Llama-3-8B produced by STAMP and STAMP-LR at their default configurations ($r = 64$, 10\% retain buffer). Leakage is meaningful only relative to what the model knew, so \textit{Base} provides the reference upper bound; the \textit{Oracle} never saw $\mathcal{D}_f$, so its forget-side scores are not evaluated (N/A), consistent with Table~\ref{tab:STAMP-Comparison}.

\begin{table}[htbp]
\centering
\caption{Robustness to query reformulation on personal data erasure ($F\text{-}RL\downarrow$, 200 forget items).}
\label{tab:robustness}
\small
\renewcommand{\arraystretch}{1.2}
\begin{tabular}{l|cccc}
\Xhline{1.15pt}
\textbf{Method} & \textbf{Original} & \textbf{Paraphr.} & \textbf{Aspect} & \textbf{Indirect} \\
\Xhline{1.15pt}
Base     & 0.87 & 0.84 & 0.79 & 0.61 \\
\rowcolor{gray!15}
Oracle   & N/A & N/A & N/A & N/A \\
\textbf{STAMP}    & 0.00 & 0.00 & 0.00 & 0.00 \\
\rowcolor{gray!15}
\textbf{STAMP-LR} & 0.00 & 0.00 & 0.00 & 0.00 \\
\Xhline{1.15pt}
\end{tabular}
\end{table}

Table~\ref{tab:robustness} shows that forgetting survives reformulation. Both variants hold $F\text{-}RL$ at 0.00 across paraphrased, aspect-shifted, and indirect queries: the healed model answers every reformulation with a refusal statement, which shares no lexical content with the erased biography, so ROUGE-L collapses to zero. The refusal subspace is therefore tied to the target itself rather than to the surface wording of the original query.
\subsection{Ablation}
\noindent

\textbf{Layer, rank, and retain-buffer ablations:}
Sweeping the intervention layer leaves forgetting complete everywhere ($Acc_\textit{f} = 0.00$) and retention within 79.17--80.13, with layer~6 the best, inside the band ROME and MEMIT edit~\cite{meng2022locating, meng2022mass}. Rank ablation shows retention stable for $r \geq 64$, our default, and a 10\% retain buffer leaves retention essentially unchanged (0.88 vs.\ 0.90) while halving runtime. Full tables are in the supplementary material (Sections~C.1--C.3).




